% Discussion section. Numbers quoted here all appear in Section 5's tables,
% which are generated by cf_worldmodels/experiments/export_tables.py.

\section{Discussion}
\label{sec:discussion}

\subsection{What ``controlled dynamic distance'' can and cannot mean}
\label{sec:discussion:axis}

We built this benchmark around three levels of dynamic distance per family and
described them as controlled. The grid says that description needs amending, and
the amendment is the most portable thing here, because the same construction is
standard practice: continual-learning suites routinely order task sequences by a
qualitative notion of similarity and treat that ordering as an experimental
factor.

What is genuinely controlled is the \emph{construction}, and in Gymnasium it is
controlled in the strongest sense available: the maximum level is the medium
perturbation plus two further ones, applied to the same task pair. If the axis
were ordering anything monotone, a strict superset of a perturbation could not
come out below it. It does, by a wide margin
(Section~\ref{sec:results:axis}), and the measured distance sides with the
outcome rather than with the construction.

What is not controlled, then, is the quantity the label is taken to stand for.
The ordering of the labels survives contact with neither of the two things we
can measure: not the forgetting the cells produce, and not the measured distance
between the environments' transition models, which additionally inverts
MiniGrid's minimum and medium levels. A label that neither predicts the outcome
nor tracks the measured construct is a name, not a control.

The practical recommendation is narrow and cheap: \textbf{report a measured
distance for every task pair, and treat level labels as construction parameters
rather than as an independent variable.} $d_{\mathrm{trans}}$ costs one model per
environment --- in our grid, a 20\% surcharge on the total training budget --- and
it is defined wherever transitions can be sampled, including families where no
parameter vector exists to difference.

Our own instrument is not the finished form of that recommendation, and we
would rather say so than let $\rho = 0.58$ stand as an endorsement. Two defects,
both established in Section~\ref{sec:results:predictors}. It does not resolve
orderings within a family at five seeds: the gaps between Gymnasium's levels are
a third of the spread within them, and doubling the training budget swaps two of
the three. And as we compute it, the two reference models are trained
independently, so one is scored in the other's latent basis and part of what it
reports is basis mismatch rather than a difference in dynamics
(Section~\ref{sec:method:distances}) --- which also explains why the quantity
moves when the budget does, and a distance between two environments should not.

Both have concrete fixes we did not run: a reference pair sharing one frozen
encoder trained on both tasks removes the confound, and more seeds resolve the
orderings or show there is nothing to resolve. The recommendation we stand
behind is therefore the weaker and more durable half of it --- report a measured
distance and hold it to account --- rather than this particular estimator. It is
worth noting which way the evidence ran: it was the measured quantity that
exposed its own limits under a budget change, and the labelled axis that could
not, because a label has nothing to be wrong about.

\subsection{Two things called ``distance'' are being conflated}
\label{sec:discussion:conflation}

The peak at the medium level is consistent across families, which invites a
mechanism rather than a list of exceptions. The one we favour is that forgetting
tracks \emph{how much the model actually updates}, and that this is non-monotone
in distance for a reason that has nothing to do with similarity: overwriting
requires a task that is both different and learnable, and the two conditions
pull apart at the top of the axis.

Gymnasium is the case where we can say why, because the budget rerun measured
it. There the maximum level is not a task the model fails to learn --- it fits
task~B \emph{better} than the medium level does, at both budgets. A cheetah at
triple mass, half friction and $40\%$ gravity moves very little, and the
observations it generates are correspondingly less varied. The strongest
perturbation in that family produces the \emph{easiest} task~B, and the least
forgetting, which is the account working rather than an exception to it: what
was overwritten tracked what there was to learn, not how far the environment had
been moved. Note also what this does to the axis, since a perturbation is
supposed to make a task harder: the construction's own premise fails at its own
top level.

If that is right, the notion of task distance used in this literature bundles
two variables that come apart. One is \emph{dissimilarity}: how far the new
environment's dynamics sit from the old. The other is \emph{learning demand}: how
much the model has to change to fit the new task at the available budget. They
correlate at the low end --- a nearby task is easy and demands little --- and they
decouple at the high end, where a distant task can be one the model learns
almost nothing from: because it is too hard, as we suspect in DMControl, or
because it is too impoverished, as we measured in Gymnasium. Those are opposite
failures of the same bundling.

Our data show the seam. Task-B difficulty is the best single predictor we have
($\rho = 0.58$), and it is best precisely because it is closer to learning
demand than to dissimilarity. But a single scalar cannot capture both arms of a
non-monotone relation, which is why the correlation is $0.58$ and not higher, and
why DMControl breaks the account: there task~B gets steadily harder to fit from
minimum to maximum ($15.43$, $34.07$, $40.54$) while forgetting falls after the
medium level --- the regime where difficulty has stopped meaning ``much to learn''
and started meaning ``too hard to learn''. Separating the two would need a
predictor built from what the optimiser did during task~B rather than from the
outcome, and we did not run that experiment.

\subsection{Forgetting is not where the metrics were looking}
\label{sec:discussion:encoder}

The second result is more actionable than the first, and it is uncomfortable for
the metrics we ourselves proposed.

PF and RD are computed on latents encoded once by the post-task-A model. That is
what isolates the transition component, which is the benchmark's stated purpose,
and it is the right choice for that purpose: letting each model encode with its
own drifting VAE would score each against its own moving coordinates rather than
against a common evaluation set. The consequence is that both metrics are blind
to drift in the encoder, and the blindness is not a rounding error. Fine-tuning's
held-out task-A reconstruction grows by a factor of $811$ in MiniGrid at minimum
distance while its PF reads $-5.58$ --- in the frozen basis, an improvement. Its
PF is negative in seven of nine cells.

The lesson generalises past our particular metrics: \textbf{any latent-space
forgetting metric evaluated on frozen representations measures the component it
was pointed at and reports nothing about the representation itself}, and in an
end-to-end world model the representation is where most of the damage is. A
benchmark that reports only such metrics will report that a fine-tuned model did
not forget. Recording a representation-space quantity alongside them costs one
forward pass per evaluation and turns a silent scope limitation into a reported
number.

EWC makes the same point structurally, and it is the cleanest result in the
paper because it was derived before it was measured. EWC's Fisher information is
defined over $\log P(z' \mid z, a)$; the encoder parameters do not appear in that
term, so their Fisher entries are \emph{exactly} zero --- not small, zero --- and
the quadratic penalty cannot constrain them even in principle. The prediction has
two halves and both hold: EWC preserves the transition component almost perfectly
($|\text{PF}| \le 0.07$ in seven of nine cells) while its encoder degrades by a
factor of $800$, indistinguishable from the unregularised baseline's $811$. The Fisher is
also exactly zero on \texttt{gru.weight\_hh}, because the Fisher set consists of
single transitions started from $h = 0$, so the recurrent pathway is unprotected
as well.

This is worth stating as a general caution about regularisation-based methods:
\textbf{the protected set is exactly the support of whatever likelihood the
Fisher is estimated from}, and in modular architectures that support is usually
narrower than the module list suggests. The failure is invisible to any metric
sharing the same restriction --- which, in our first formulation, was every
metric we reported.

\subsection{What the three mitigation families actually buy}
\label{sec:discussion:methods}

Read across both scales, the three standard strategies separate cleanly, and none
of them dominates.

\textbf{Replay} is the only method that keeps both components: its task-A
reconstruction never grows by more than $31\%$ and improves outright in
DMControl, and it has the lower RD in four of the six cells that produce
forgetting, unanimously across ten seeds in each. It pays for this where it
works hardest. In the three non-MiniGrid cells where it protects the encoder,
its forward transfer is the worst of the four RSSM methods by a factor of three
or more ($-8.52$, $-8.06$ and $-7.40$, against next-worst values of $-2.15$,
$-1.85$ and $-2.36$). Part
of that is an artefact we must declare: replay trains on A and B during the
task-B phase, so half its gradient steps go to task-A data and its FT confounds
transfer with a halved effective budget on the new task. The rest is the
expected price of rehearsal, and it is a price the metric can see only because FT
is measured against a model trained on task~B from scratch.

\textbf{Regularisation} protects what its penalty covers and nothing else, as
above. Its appeal --- no stored data, no growing parameter count --- is intact; what
our grid shows is that its guarantees are component-scoped in a way that a
system-level evaluation cannot detect.

\textbf{Structural isolation}, including our own method, buys preservation by
declining to learn. UG-MTM's task-A reconstruction is unchanged to the bit in all
nine cells, because its encoder is frozen; its forward transfer is $-557$, $-526$
and $-1642$ in MiniGrid, where no other method leaves $[-8.52, +2.01]$ anywhere
in the grid. The honest summary is that it does not forget because it does not
learn. The cost nearly vanishes when the two tasks look alike ($-4.33$ in
Gymnasium at minimum distance), which locates a regime where freezing is
affordable rather than showing that the mechanism works. We report this because
the benchmark produced it; a suite whose author's method came out ahead on the
author's own benchmark would be worth less.

Freezing has a second consequence we did not anticipate and would not have found
without the seed-level view. UG-MTM's post-task-B transition models become
\emph{overconfident}: in the extreme seed of MiniGrid's maximum-distance cell,
some latent dimensions collapse to $\sigma = 0.007$, the rollout KL is already
$80$ at step~0, and RD reaches $4364$ on that seed against a median of $576$. Confidence
without accuracy is a failure mode that an unbounded divergence will surface and
a bounded one will hide, and it is invisible to any aggregate reported as a mean.

\subsection{Lessons for building this kind of benchmark}
\label{sec:discussion:measurement}

Three of our design decisions were forced by measurements rather than chosen in
advance, and we think they generalise.

\textbf{Report medians when the metric is unbounded.} RD is a KL divergence, so
methods that produce overconfident models get heavy right tails by construction.
This is not rare in our grid: 17 of 180 forgetting cells are right-skewed at
five seeds. Means over five seeds describe those cells badly, and the tail is
itself a result about the method rather than noise to be smoothed, so we report
medians with the full seed range and mark the skewed cells.

\textbf{Declare the cells that measure nothing.} Three of our nine cells produce
no forgetting in any method. That is the correct behaviour for the bottom of a
distance axis and evidence that the axis reaches down far enough, but it also
means those cells discriminate between methods on nothing, and reporting them
without saying so inflates the apparent size of the grid by half. Our effective
grid is six cells.

\textbf{Aggregate with an explicit inclusion rule.} The family-level numbers in
Section~\ref{sec:results:axis} exclude UG-MTM, whose frozen encoder puts its RD
two orders of magnitude below the others in DMControl. That exclusion changes the
headline numbers, so it belongs in the paper and in the code that produces the
tables rather than in an analysis script nobody sees.

\subsection{Limitations, and what would change these conclusions}
\label{sec:discussion:limitations}

Statistical resolution was the binding constraint, and we spent compute on it
rather than argument. The six cells that discriminate between methods run ten
seeds, which lowers the exact paired test's floor from $0.0625$ to about
$0.002$ and is why the comparisons in Section~\ref{sec:discussion:methods} can
be significant at all. It also changed one of them: at five seeds replay looked
better than fine-tuning in five of six forgetting cells, and at ten the sixth
turns out to be a wash ($p = 0.68$) rather than a win. That is the resolution
doing its job, and it is a reminder that a five-seed direction is a direction,
not a result. Ten is still not many, and the three control cells remain at
five; more seeds is still the cheapest improvement available.

The grid is nine cells of which six discriminate, one task switch rather than a
sequence, and one architecture family: every method here is a variation on the
same recurrent latent world model, so the results speak to that class and not to
world models in general. Our own measured distance inherits a confound we
declare but did not remove --- it scores one environment's transition model in
the other's latent basis --- which bounds how much weight the
$d_{\mathrm{trans}}$ half of Section~\ref{sec:results:predictors} can carry. The
budget --- 20 random-policy episodes per task and
5000 gradient updates --- is small enough that absolute magnitudes are
budget-dependent, although task~A is demonstrably learned at this scale, which is
the precondition a forgetting benchmark has to establish before any of its
numbers mean anything.

Two specific results would be worth attacking. The peak at the medium level is
the finding we are most confident in, since it repeats in three of three families
and survives changes of aggregation, metric and threshold; it would be
substantially strengthened by a family in which the levels are constructed to be
monotone in $d_{\mathrm{trans}}$ and the peak persists anyway. We did not do this
here, and we note that it has to be done carefully: choosing levels so that the
measured distance comes out monotone, and then reporting that forgetting follows
it, would be circular. The prediction that distinguishes our account from the
design's is testable and cheap --- forgetting should fall again at distances large
enough that task~B stops being learnable at the given budget.

The second is policy-level impact. An earlier formulation of this benchmark
announced a policy-impact metric alongside the three we report; it was never
implemented, because measuring it means training a controller inside the model's
imagination and evaluating it in the real environment, and we withdraw it rather
than report a placeholder. It remains the most interesting gap: everything here
measures degradation of the model, and whether the ordering we find survives the
translation into control performance is an open question that this suite, as it
stands, cannot answer.

\subsection{On the previous version of these results}
\label{sec:discussion:previous}

An earlier version of this benchmark reported five findings, at significance
levels a five-seed design cannot produce. None of them survives re-measurement,
and the causes were defects in the instrument rather than in the analysis: a
collapsed VAE posterior that fed the transition model a constant latent, an
evaluation set of Gaussian noise, a next-state objective scored against the wrong
target, a mis-specified Gaussian KL that inflated RD roughly twelvefold, and
unseeded environments that made five seeds five draws of different data rather
than five replicates.

We report this deliberately. The strongest of those findings --- that replay is
insufficient in low-capacity settings --- was the one that reversed, and it was
also the most counterintuitive, which is the pattern to expect when a measurement
is broken: an artefact is more likely to look surprising than a real effect is.
The infrastructure that replaced it is the substrate the rest of this paper rests
on. Runs reproduce bit for bit across processes and across days: a cell
re-executed eight days after the original run, to diagnose the heavy tail of
Section~\ref{sec:discussion:methods}, returned $\text{RD} = 4364.4502$ against
$4364.4502$ stored and $\text{PF} = 39.9796$ against $39.9796$. Every result
records the protocol it ran under and the task pair it came from, and the runner
refuses to combine results produced under different ones. Those properties are
unglamorous and they are what makes the negative result in
Section~\ref{sec:results:axis} worth reporting at all.
